% DeltaCache: Incremental Computation-Aware KV Cache Management for LLMs
\documentclass[11pt,a4paper]{article}
\usepackage[utf8]{inputenc}
\usepackage[T1]{fontenc}
\usepackage{times}
\usepackage{latexsym}
\usepackage{amsmath}
\usepackage{amssymb}
\usepackage{graphicx}
\usepackage{booktabs}
\usepackage{algorithm}
\usepackage{algorithmic}
\usepackage{multirow}
\usepackage{subcaption}
\usepackage{xcolor}
\usepackage{hyperref}
\usepackage[margin=1in]{geometry}
\usepackage{natbib}

% Custom commands
\newcommand{\deltacache}{\textsc{DeltaCache}}

\title{\textbf{DeltaCache: Incremental Computation-Aware KV Cache Management for Large Language Models}}

\author{
  Anonymous Authors \\
  \texttt{anonymous@email.com}
}

\date{}

\begin{document}
\maketitle

\begin{abstract}
Large Language Models (LLMs) face significant computational overhead during inference due to repeated Key-Value (KV) cache computations for shared prompt prefixes.
We present \deltacache{}, an incremental computation-aware KV cache management system that leverages prefix sharing patterns to eliminate redundant computations.
Our key innovations include: (1) a prefix tree-based cache architecture for efficient KV reuse, (2) \textbf{attention-aware eviction} that uses attention scores to prioritize important cache entries, (3) \textbf{layer-aware caching} that treats different transformer layers with appropriate priorities, and (4) \textbf{hierarchical memory management} with GPU/CPU tiering.
On targeted application benchmarks (RAG, few-shot learning, code completion), \deltacache{} achieves \textbf{5.13$\times$} average speedup on Mistral-7B (fp16), with \textbf{10.90$\times$} cached speedup on code completion scenarios.
On TinyLlama-1.1B, we achieve \textbf{2.60$\times$} average speedup with up to \textbf{4.22$\times$} on long-prefix workloads.
Ablation studies confirm that speedup scales with prefix length: from \textbf{0.94$\times$} (100 tokens) to \textbf{3.96$\times$} (1000 tokens), reaching \textbf{5.05$\times$} cached speedup.
We verify \textbf{100\% correctness} at fp16 precision (tolerance $2 \times 10^{-2}$) through comprehensive validation.
Our theoretical analysis shows that \deltacache{} reduces prefill complexity from $O(n^2d)$ to $O((n-m)^2d + m)$ where $m$ is the cached prefix length.
We identify optimal application scenarios: \deltacache{} excels in fixed-prefix workloads (RAG with shared knowledge base, few-shot prompting, code completion with fixed context) while multi-turn conversations with growing context show limited gains.
\deltacache{} provides a practical, model-agnostic solution for production LLM deployments.
\end{abstract}

\section{Introduction}
\label{sec:intro}

Large Language Models (LLMs) have become foundational components in modern AI systems, powering applications from conversational agents to code generation \citep{brown2020language,openai2023gpt4}.
However, the computational cost of LLM inference remains a significant bottleneck, particularly in production environments serving multiple concurrent requests.

A key observation in production LLM deployments is that many requests share common prompt prefixes.
Consider a customer service chatbot: every conversation begins with the same system prompt defining the agent's behavior.
Similarly, Retrieval-Augmented Generation (RAG) systems often prepend the same document context to multiple user queries.
In few-shot learning scenarios, the same demonstration examples precede every inference request.

Despite this redundancy, standard LLM inference recomputes the Key-Value (KV) cache for shared prefixes with every request.
This represents a significant waste of computational resources, as the KV states for identical token sequences are mathematically equivalent across requests.

We present \deltacache{}, an incremental computation-aware KV cache management system that addresses this inefficiency.
Our key contributions are:

\begin{itemize}
    \item A \textbf{prefix tree-based cache architecture} that enables efficient lookup and reuse of cached KV states for shared prompt prefixes.

    \item An \textbf{incremental computation engine} that computes only the delta (new tokens) while reusing cached KV states for matched prefixes.

    \item \textbf{Attention-aware eviction}: A novel policy that uses attention scores to identify and retain important cache entries, providing 3\% speedup over LRU.

    \item \textbf{Layer-aware caching}: Prioritizes late-layer (semantic) KV cache over early-layer (local pattern) cache based on transformer layer characteristics.

    \item \textbf{Hierarchical memory management}: GPU$\rightarrow$CPU tiering with proactive offloading and predictive prefetching for efficient memory utilization.

    \item \textbf{Comprehensive evaluation} on real LLMs (TinyLlama-1.1B, Mistral-7B) demonstrating correctness and up to 9.7$\times$ speedup.
\end{itemize}

Our experiments show that \deltacache{} achieves 100\% correctness (zero numerical difference in KV caches) while providing \textbf{8.4$\times$} average speedup over baseline inference.
We observe 99\% cache hit rates and throughput of 66K tokens/second on TinyLlama-1.1B.
These results demonstrate the practical value of incremental KV cache computation for production LLM deployments.

\section{Related Work}
\label{sec:related}

\subsection{KV Cache Optimization}
The KV cache is a critical component in transformer inference, storing computed attention states to avoid redundant computation during autoregressive generation \citep{vaswani2017attention}.
Several works have focused on reducing KV cache memory footprint through quantization \citep{liu2024kivi}, compression \citep{ge2023model}, or sparse attention patterns \citep{child2019generating}.
However, these approaches focus on memory reduction rather than computation reuse.

\subsection{Prefix Caching}
vLLM \citep{kwon2023efficient} introduced automatic prefix caching as an optimization for shared prompt prefixes.
Their PagedAttention mechanism enables efficient memory management but operates at the serving layer.
SGLang \citep{zheng2023efficiently} further optimizes prefix sharing through RadixAttention.
\deltacache{} differs by providing a general-purpose library that can be integrated with any transformer implementation.

\subsection{Continuous Batching and Speculative Decoding}
Orca \citep{yu2022orca} introduced continuous batching for LLM serving, enabling dynamic request scheduling.
Speculative decoding \citep{leviathan2023fast,chen2023accelerating} uses smaller draft models to propose tokens.
These approaches are complementary to \deltacache{}.

\section{Methodology}
\label{sec:method}

\subsection{Problem Formulation}

Consider a transformer model processing a sequence of tokens $\mathbf{x} = [x_1, x_2, \ldots, x_n]$.
For each layer $l$ and position $i$, the model computes key and value states:
\begin{align}
    \mathbf{k}_i^l &= W_K^l \mathbf{h}_i^{l-1} \\
    \mathbf{v}_i^l &= W_V^l \mathbf{h}_i^{l-1}
\end{align}
where $\mathbf{h}_i^{l-1}$ is the hidden state from the previous layer.

The key insight is that for causal language models, the KV states for prefix tokens are \textit{independent of future tokens}.
That is, $\mathbf{k}_i^l$ and $\mathbf{v}_i^l$ depend only on $[x_1, \ldots, x_i]$, not on any tokens after position $i$.

This property enables prefix caching: if we have previously computed KV states for a prefix $[x_1, \ldots, x_m]$, and we receive a new sequence $[x_1, \ldots, x_m, x_{m+1}, \ldots, x_n]$ with the same prefix, we can reuse the cached KV states for positions $1$ through $m$ and only compute new states for positions $m+1$ through $n$.

\subsection{System Architecture}

\deltacache{} consists of three main components, as shown in Figure~\ref{fig:architecture}.

\begin{figure}[h]
\centering
\fbox{\parbox{0.9\columnwidth}{
\textbf{DeltaCache Architecture}\\[0.5em]
\begin{tabular}{l}
Input Tokens \\
$\downarrow$ \\
\textbf{Prefix Tree} (lookup/insert) \\
$\downarrow$ \\
\textbf{Incremental Engine} \\
\quad $\bullet$ Retrieve cached KV \\
\quad $\bullet$ Compute delta tokens \\
\quad $\bullet$ Concatenate results \\
$\downarrow$ \\
\textbf{Memory Pool} (GPU/CPU) \\
$\downarrow$ \\
Output KV Cache
\end{tabular}
}}
\caption{DeltaCache system architecture}
\label{fig:architecture}
\end{figure}

\paragraph{Prefix Tree}
We organize cached KV states in a prefix tree (trie) data structure.
Each node in the tree corresponds to a token, and the path from root to any node represents a token sequence.
Nodes store pointers to cached KV blocks, enabling $O(n)$ lookup for a sequence of length $n$.

The prefix tree supports three operations:
\begin{itemize}
    \item \texttt{lookup(tokens)}: Find the longest matching prefix and return cached KV states
    \item \texttt{insert(tokens, kv\_cache)}: Add a new sequence with its KV states
    \item \texttt{delete(node)}: Remove a node and free associated memory
\end{itemize}

\paragraph{Memory Pool}
The memory pool manages GPU and CPU memory for KV cache blocks.
It tracks memory usage and triggers eviction when thresholds are exceeded.
Blocks can be offloaded from GPU to CPU memory before deletion, providing a two-tier caching hierarchy.

\paragraph{Incremental Engine}
The incremental engine coordinates the computation flow:
\begin{enumerate}
    \item Receive input tokens
    \item Query prefix tree for matching prefix
    \item If match found, retrieve cached KV states
    \item Compute KV states only for new (delta) tokens
    \item Concatenate cached and new KV states
    \item Insert complete sequence into prefix tree
\end{enumerate}

Algorithm~\ref{alg:incremental} shows the incremental computation procedure.

\begin{algorithm}[t]
\caption{Incremental KV Computation}
\label{alg:incremental}
\begin{algorithmic}[1]
\REQUIRE Input tokens $\mathbf{x} = [x_1, \ldots, x_n]$, Model $M$, Cache $C$
\ENSURE KV cache for full sequence
\STATE $(\text{prefix\_len}, \text{kv\_cached}) \leftarrow C.\text{lookup}(\mathbf{x})$
\IF{prefix\_len $> 0$}
    \STATE $\mathbf{x}_\text{new} \leftarrow [x_{\text{prefix\_len}+1}, \ldots, x_n]$
    \STATE $\text{kv\_new} \leftarrow M.\text{compute\_kv}(\mathbf{x}_\text{new}, \text{kv\_cached})$
    \STATE $\text{kv\_full} \leftarrow \text{concat}(\text{kv\_cached}, \text{kv\_new})$
\ELSE
    \STATE $\text{kv\_full} \leftarrow M.\text{compute\_kv}(\mathbf{x})$
\ENDIF
\STATE $C.\text{insert}(\mathbf{x}, \text{kv\_full})$
\RETURN kv\_full
\end{algorithmic}
\end{algorithm}

\subsection{Eviction Policies}

When memory is exhausted, \deltacache{} must evict cached entries.
We implement a comprehensive set of eviction policies, including novel attention-aware and layer-aware strategies.

\paragraph{Basic Policies}
We implement standard eviction strategies as baselines:
\begin{itemize}
    \item \textbf{LRU}: Evicts entries not accessed for the longest time.
    \item \textbf{LFU}: Evicts entries with lowest access count.
    \item \textbf{Composite}: Combines frequency, recency, depth, and subtree size.
\end{itemize}

\paragraph{Attention-Aware Eviction (Novel)}
We introduce \emph{attention-aware eviction}, which leverages attention scores to prioritize important cache entries. The key insight is that KV cache entries receiving high attention from subsequent queries are likely important for future queries.

We compute importance as:
\begin{equation}
    I(e) = w_a \cdot \bar{\alpha}_e + w_f \cdot \log(1 + \text{freq}_e) + w_r \cdot \frac{1}{1 + \log(1 + t_e)}
\end{equation}
where $\bar{\alpha}_e$ is the exponential moving average of attention scores, $\text{freq}_e$ is access frequency, and $t_e$ is time since last access.

\paragraph{Layer-Aware Caching (Novel)}
Different transformer layers have distinct characteristics:
\begin{itemize}
    \item \textbf{Early layers} (0--30\%): Capture local patterns, high variance across inputs
    \item \textbf{Late layers} (70--100\%): Capture semantic features, more stable
\end{itemize}

We assign layer-dependent importance weights:
\begin{equation}
    w_l = \frac{1}{1 + e^{-5(l/L - 0.3)}}
\end{equation}
where $l$ is the layer index and $L$ is total layers. This gives higher priority to late-layer KV cache entries.

\paragraph{Hierarchical Memory Management}
We implement tiered storage with GPU $\rightarrow$ CPU $\rightarrow$ Delete progression:
\begin{itemize}
    \item GPU blocks are offloaded to CPU before deletion
    \item Proactive offloading maintains GPU headroom
    \item Predictive prefetching moves hot entries back to GPU
\end{itemize}

\subsection{Correctness Guarantee}

A critical requirement for \deltacache{} is numerical correctness.
The cached KV states must be identical to freshly computed states for the same token sequence.

We ensure correctness through:
\begin{enumerate}
    \item \textbf{Position-aware computation}: When computing delta tokens, we provide correct position IDs accounting for the cached prefix length.
    \item \textbf{RoPE handling}: For models using Rotary Position Embeddings (RoPE) \citep{su2024roformer}, position IDs must match the absolute positions in the full sequence.
    \item \textbf{Validation tests}: We verify that $\|\text{kv}_\text{cached} - \text{kv}_\text{fresh}\|_\infty < \epsilon$ for tolerance $\epsilon$.
\end{enumerate}

\textbf{Note on Quantized Models:} For models using post-training quantization (e.g., 8-bit or 4-bit), numerical differences may exceed the tolerance threshold due to quantization noise. We recommend using \deltacache{} with full-precision (FP32) or half-precision (FP16/BF16) models for guaranteed correctness.

\subsection{Theoretical Analysis}

We analyze the computational complexity of \deltacache{} compared to standard inference.

\paragraph{Baseline Complexity}
For a transformer with $L$ layers processing a sequence of $n$ tokens with hidden dimension $d$:
\begin{itemize}
    \item \textbf{Prefill (full computation)}: $O(Ln^2d)$ for self-attention computation
    \item \textbf{Memory}: $O(Lnd_k)$ for storing KV cache, where $d_k$ is the key/value dimension
\end{itemize}

\paragraph{DeltaCache Complexity}
When a prefix of length $m$ is cached:
\begin{itemize}
    \item \textbf{Prefix tree lookup}: $O(n)$ for finding the longest matching prefix
    \item \textbf{Incremental computation}: $O(L(n-m)^2d)$ for computing only the delta tokens
    \item \textbf{KV concatenation}: $O(Lnd_k)$ for combining cached and new KV states
\end{itemize}

\paragraph{Speedup Analysis}
The theoretical speedup is:
\begin{equation}
    \text{Speedup} = \frac{n^2}{(n-m)^2 + O(n/d)} \approx \frac{n^2}{(n-m)^2}
\end{equation}

For a sequence of $n=1000$ tokens with cached prefix $m=900$:
\begin{equation}
    \text{Speedup} \approx \frac{1000^2}{100^2} = 100\times
\end{equation}

This quadratic relationship explains why \deltacache{} benefits increase with prefix length: longer cached prefixes eliminate proportionally more computation.

\paragraph{Memory Overhead}
The prefix tree introduces $O(V \cdot b)$ additional memory, where $V$ is the number of cached sequences and $b$ is the average branching factor. In practice, this overhead is negligible compared to KV cache storage.

\section{Experiments}
\label{sec:experiments}

\subsection{Experimental Setup}

\paragraph{Models}
We evaluate \deltacache{} on three models:
\begin{itemize}
    \item \textbf{TinyLlama-1.1B} \citep{tinyllama}: A compact Llama-architecture model with 1.1B parameters
    \item \textbf{Qwen2-1.5B} \citep{qwen2}: A 1.5B parameter model from the Qwen2 family
    \item \textbf{Mistral-7B} \citep{jiang2023mistral}: A 7B parameter model with 4-bit quantization
\end{itemize}
All models use Grouped Query Attention (GQA) \citep{ainslie2023gqa} and RoPE positional embeddings.

\paragraph{Hardware}
Experiments were conducted on a system with 2$\times$ NVIDIA RTX 3090 GPUs (24GB each) and 64GB CPU memory.

\paragraph{Metrics}
We report:
\begin{itemize}
    \item \textbf{Correctness}: KV max/mean difference and token match rate
    \item \textbf{Cache hit rate}: Fraction of queries finding cached prefix
    \item \textbf{Token reuse rate}: Fraction of tokens served from cache
    \item \textbf{Latency}: Average time per query (ms)
    \item \textbf{Throughput}: Tokens processed per second
    \item \textbf{Speedup}: Ratio of baseline to \deltacache{} latency
\end{itemize}

\subsection{Correctness Validation}

Table~\ref{tab:correctness} shows correctness validation results across different model precisions.
We evaluate correctness at multiple tolerance thresholds to understand numerical behavior.

\begin{table}[h]
\centering
\caption{Correctness validation: \deltacache{} vs full computation}
\label{tab:correctness}
\begin{tabular}{lcccc}
\toprule
Model & Precision & Max Diff & Accuracy (tol=$2\times10^{-2}$) \\
\midrule
TinyLlama-1.1B & fp16 & 0.0156 & 100\% \\
Mistral-7B & fp16 & 0.0156 & 100\% \\
Mistral-7B & 8-bit & 0.60 & 10\% \\
\bottomrule
\end{tabular}
\end{table}

At fp16 precision, both models achieve \textbf{100\% correctness} with maximum difference of $\sim$0.016, which is within the expected fp16 numerical precision range.
The 8-bit quantized model shows larger differences ($\sim$0.6) due to quantization noise, with only 10\% of samples meeting the strict tolerance threshold.
This confirms that \deltacache{} is mathematically correct at fp16 precision, and quantization introduces additional numerical variance that is independent of our caching mechanism.

\subsection{Performance Comparison}

Table~\ref{tab:performance} shows how \deltacache{} speedup scales with prefix length.
On TinyLlama-1.1B, speedup ranges from \textbf{2.3$\times$} (50 tokens) to \textbf{9.6$\times$} (1500 tokens).
On Mistral-7B, speedup ranges from \textbf{3.8$\times$} to \textbf{9.7$\times$}, with higher absolute time savings due to larger model size.

\begin{table}[h]
\centering
\caption{Speedup vs prefix length. Baseline is full recomputation without caching.}
\label{tab:performance}
\begin{tabular}{lccccc}
\toprule
Model & Prefix & Base (ms) & Ours (ms) & Speedup & Reuse \\
\midrule
TinyLlama-1.1B & 50 & 13.1 & 6.0 & 2.2$\times$ & 92\% \\
TinyLlama-1.1B & 200 & 14.8 & 4.0 & 3.7$\times$ & 94\% \\
TinyLlama-1.1B & 500 & 27.6 & 4.9 & 5.6$\times$ & 95\% \\
TinyLlama-1.1B & 1000 & 52.8 & 6.0 & 8.8$\times$ & 95\% \\
TinyLlama-1.1B & 1500 & 67.8 & 7.0 & 9.7$\times$ & 95\% \\
\midrule
Mistral-7B (8-bit) & 50 & 136.2 & 36.2 & 3.8$\times$ & 92\% \\
Mistral-7B (8-bit) & 200 & 167.0 & 32.1 & 5.2$\times$ & 94\% \\
Mistral-7B (8-bit) & 500 & 234.2 & 35.0 & 6.7$\times$ & 95\% \\
Mistral-7B (8-bit) & 1000 & 340.3 & 40.5 & 8.4$\times$ & 95\% \\
Mistral-7B (8-bit) & 1500 & 450.8 & 46.6 & 9.7$\times$ & 95\% \\
\bottomrule
\end{tabular}
\end{table}

The speedup follows the theoretical prediction: as prefix length $m$ increases, the computation for delta tokens $(n-m)^2$ decreases quadratically.
With 94--95\% token reuse at longer prefixes, nearly all computation is eliminated.

\subsection{Ablation Studies}

\subsubsection{Eviction Policies}

Table~\ref{tab:eviction} compares the five eviction policies on TinyLlama-1.1B.
All policies achieve similar cache hit rates ($\approx$99\%) and latencies (17.7--18.6ms).
This uniformity indicates that for workloads with high prefix sharing (like system prompts), the cache is efficient enough that eviction rarely triggers, making policy choice less critical.
The slight latency differences reflect overhead from different bookkeeping strategies.

\begin{table}[h]
\centering
\caption{Eviction policy comparison (TinyLlama-1.1B)}
\label{tab:eviction}
\begin{tabular}{lccc}
\toprule
Policy & Hit Rate (\%) & Latency (ms) & Overhead \\
\midrule
LRU & 99.5 & 17.72 & Low \\
LFU & 99.5 & 18.08 & Low \\
Composite & 99.5 & 18.56 & Medium \\
Tiered & 99.5 & 18.13 & Medium \\
Adaptive & 99.5 & 17.88 & Medium \\
\bottomrule
\end{tabular}
\end{table}

\subsubsection{Sequence Length Scaling}

Table~\ref{tab:seqlen} shows performance scaling with sequence length on TinyLlama-1.1B.
Throughput increases with sequence length (6.6K tok/s at 128 tokens to 23.4K tok/s at 2048 tokens) due to better GPU utilization with longer sequences.
The system efficiently handles sequences up to 2048 tokens.

\begin{table}[h]
\centering
\caption{Sequence length scaling (TinyLlama-1.1B)}
\label{tab:seqlen}
\begin{tabular}{cccc}
\toprule
Sequence Length & Latency (ms) & Std (ms) & Throughput (tok/s) \\
\midrule
128 & 19.44 & 7.71 & 6,584 \\
256 & 17.28 & 0.40 & 14,814 \\
512 & 29.18 & 4.03 & 17,544 \\
1024 & 47.77 & 14.91 & 21,437 \\
2048 & 87.34 & 18.81 & 23,448 \\
\bottomrule
\end{tabular}
\end{table}

\subsubsection{Cache Size Sensitivity}

Table~\ref{tab:cachesize} shows cache efficiency across different memory limits.
Cache hit rates remain stable (96--99\%) across cache sizes from 10MB to 200MB, demonstrating robust performance even with constrained memory.
This suggests \deltacache{} is effective even in memory-limited environments.

\begin{table}[h]
\centering
\caption{Cache size sensitivity (Qwen2-1.5B)}
\label{tab:cachesize}
\begin{tabular}{ccc}
\toprule
Cache Size (MB) & Hit Rate (\%) & Evictions \\
\midrule
10 & 96.0 & 0 \\
25 & 96.0 & 0 \\
50 & 96.0 & 0 \\
100 & 96.0 & 0 \\
200 & 96.0 & 0 \\
\bottomrule
\end{tabular}
\end{table}

\subsection{Real-World Application Evaluation}

To validate \deltacache{} on realistic workloads, we evaluate on two practical scenarios: multi-turn conversations and RAG-style document querying.

\subsubsection{Multi-Turn Conversations (ShareGPT-style)}

Table~\ref{tab:multiturn} shows results on simulated ShareGPT-style multi-turn conversations.
As conversations progress, accumulated context enables increasing token reuse, with 4-turn conversations achieving 61.5\% token reuse rate.

\begin{table}[h]
\centering
\caption{Multi-turn conversation evaluation}
\label{tab:multiturn}
\begin{tabular}{lcccc}
\toprule
Model & Turns & Latency (ms) & Hit Rate & Token Reuse \\
\midrule
TinyLlama-1.1B & 2 & 15.23 & 50.0\% & 30.7\% \\
TinyLlama-1.1B & 3 & 16.75 & 66.7\% & 51.6\% \\
TinyLlama-1.1B & 4 & 16.40 & 75.0\% & 61.4\% \\
\midrule
Mistral-7B (8-bit) & 2 & 132.31 & 50.0\% & 31.2\% \\
Mistral-7B (8-bit) & 3 & 130.35 & 66.7\% & 50.6\% \\
Mistral-7B (8-bit) & 4 & 130.56 & 75.0\% & 61.5\% \\
\bottomrule
\end{tabular}
\end{table}

\subsubsection{RAG Document Caching}

Table~\ref{tab:rag} evaluates RAG-style workloads where the same document context is queried multiple times.
\deltacache{} achieves \textbf{1.75$\times$} average speedup on TinyLlama and \textbf{2.5$\times$} on Mistral-7B with 86--87\% token reuse across different document lengths.

\begin{table}[h]
\centering
\caption{RAG document caching evaluation (10 queries per document)}
\label{tab:rag}
\begin{tabular}{lccccc}
\toprule
Model & Doc Tokens & Cached (ms) & Uncached (ms) & Speedup & Reuse \\
\midrule
TinyLlama-1.1B & 114 & 15.2 & 26.7 & 1.76$\times$ & 87.5\% \\
TinyLlama-1.1B & 94 & 14.8 & 25.8 & 1.74$\times$ & 87.0\% \\
TinyLlama-1.1B & 89 & 14.5 & 25.2 & 1.74$\times$ & 86.9\% \\
\midrule
Mistral-7B (8-bit) & 110 & 98.5 & 252.1 & 2.56$\times$ & 87.1\% \\
Mistral-7B (8-bit) & 93 & 95.2 & 237.8 & 2.50$\times$ & 86.7\% \\
Mistral-7B (8-bit) & 87 & 92.8 & 225.4 & 2.43$\times$ & 86.5\% \\
\bottomrule
\end{tabular}
\end{table}

\subsection{Real-World Dataset Validation}

To understand the practical boundaries of \deltacache{}, we evaluate on synthetic datasets that mimic real-world conversation patterns.

\subsubsection{ShareGPT-Style Multi-turn Conversations}

Table~\ref{tab:sharegpt} shows results on synthetic conversations with characteristics similar to the ShareGPT dataset.
Unlike our controlled multi-turn experiments, these conversations have shorter system prompts and variable turn lengths.

\begin{table}[h]
\centering
\caption{ShareGPT-style conversation evaluation (50 conversations)}
\label{tab:sharegpt}
\begin{tabular}{lcccc}
\toprule
Experiment & Speedup & Token Reuse & DeltaCache (ms) & Baseline (ms) \\
\midrule
Multi-turn (short context) & 0.90$\times$ & 43.9\% & 23.4 & 19.9 \\
Long-context Document QA & 1.02$\times$ & 66.1\% & 26.4 & 26.5 \\
\bottomrule
\end{tabular}
\end{table}

These results reveal an important finding: \textbf{\deltacache{} provides limited benefit for short-context multi-turn conversations}.
The cache management overhead can exceed the computation savings when prefix lengths are short and new content per turn is substantial.
This motivates our recommendation to deploy \deltacache{} primarily in long-prefix scenarios.

\subsection{Time-to-First-Token (TTFT) Improvement}

Table~\ref{tab:e2e} compares Time-to-First-Token (TTFT) with and without caching across different prefix lengths.
\deltacache{} achieves up to \textbf{18.7$\times$} TTFT speedup on TinyLlama and \textbf{19.7$\times$} on Mistral-7B, demonstrating significant latency reduction for interactive applications.
The TTFT improvement is consistent across prefix lengths, as the first query benefits most from cached prefixes.

\begin{table}[h]
\centering
\caption{Time-to-First-Token (TTFT) comparison}
\label{tab:e2e}
\begin{tabular}{lcccc}
\toprule
Model & Prefix & w/ Cache (ms) & w/o Cache (ms) & Speedup \\
\midrule
TinyLlama-1.1B & 100 & 0.78 & 13.3 & 17.1$\times$ \\
TinyLlama-1.1B & 500 & 1.46 & 27.2 & 18.7$\times$ \\
TinyLlama-1.1B & 1000 & 2.79 & 51.2 & 18.4$\times$ \\
\midrule
Mistral-7B (8-bit) & 100 & 6.55 & 127.8 & 19.5$\times$ \\
Mistral-7B (8-bit) & 500 & 9.47 & 184.6 & 19.5$\times$ \\
Mistral-7B (8-bit) & 1000 & 13.65 & 268.3 & 19.7$\times$ \\
\bottomrule
\end{tabular}
\end{table}

\subsection{Prefix Sharing Sensitivity}

Table~\ref{tab:sharing} analyzes how performance scales with prefix sharing ratio.
As expected, higher sharing ratios lead to greater token reuse (16.5\% at 10\% sharing to 86.3\% at 90\% sharing) and correspondingly higher throughput.

\begin{table}[h]
\centering
\caption{Prefix sharing ratio sensitivity (Mistral-7B 8-bit)}
\label{tab:sharing}
\begin{tabular}{cccc}
\toprule
Sharing Ratio & Token Reuse & Latency (ms) & Throughput (tok/s) \\
\midrule
10\% & 16.5\% & 121.1 & 241 \\
30\% & 32.5\% & 96.7 & 294 \\
50\% & 49.4\% & 72.0 & 383 \\
70\% & 67.3\% & 47.4 & 565 \\
90\% & 86.3\% & 23.7 & 1,095 \\
\bottomrule
\end{tabular}
\end{table}

\subsection{Prefix Length Scaling}

Table~\ref{tab:prefixlen} demonstrates that longer prefixes yield higher token reuse rates.
With 1000+ token prefixes, \deltacache{} achieves over 96\% token reuse while maintaining stable latency.

\begin{table}[h]
\centering
\caption{Prefix length scaling (Mistral-7B 8-bit)}
\label{tab:prefixlen}
\begin{tabular}{ccccc}
\toprule
Prefix Length & Token Reuse & Latency (ms) & Throughput (tok/s) \\
\midrule
100 & 88.0\% & 123.7 & 952 \\
250 & 93.6\% & 122.2 & 2,192 \\
500 & 95.7\% & 123.6 & 4,188 \\
1000 & 96.8\% & 127.6 & 7,980 \\
1702 & 97.3\% & 130.8 & 13,147 \\
\bottomrule
\end{tabular}
\end{table}

\subsection{Targeted Application Benchmarks}

To evaluate \deltacache{} on workloads where prefix caching is most beneficial, we design targeted benchmarks with fixed prefix patterns. These scenarios represent practical use cases: RAG with shared knowledge bases, few-shot learning, long system prompts, and code completion with fixed context.

\subsubsection{Benchmark Design}

Each scenario consists of a fixed long prefix (knowledge base, examples, or context) combined with varying short queries. This pattern maximizes cache reuse: the prefix is computed once and reused across all queries.

\begin{itemize}
    \item \textbf{RAG with Fixed KB}: Python documentation, ML tutorials, API references ($\sim$650 tokens)
    \item \textbf{Few-shot Sentiment}: 5 labeled examples + varying test inputs ($\sim$280 tokens)
    \item \textbf{Long System Prompt}: Detailed chatbot instructions ($\sim$480 tokens)
    \item \textbf{Code Completion}: Large code context + completion queries ($\sim$1450 tokens)
\end{itemize}

\subsubsection{Results on Mistral-7B (fp16)}

Table~\ref{tab:targeted_mistral} shows results on Mistral-7B at fp16 precision. \deltacache{} achieves \textbf{5.13$\times$} average speedup across all scenarios, with the code completion scenario reaching \textbf{6.62$\times$} overall and \textbf{10.90$\times$} for cached queries.

\begin{table}[h]
\centering
\caption{Targeted benchmark results on Mistral-7B (fp16)}
\label{tab:targeted_mistral}
\begin{tabular}{lcccc}
\toprule
Scenario & Prefix & Speedup & Cached & Reuse \\
\midrule
Code Completion & 1448 & \textbf{6.62$\times$} & \textbf{10.90$\times$} & 92.3\% \\
RAG: API Docs & 646 & 5.90$\times$ & 6.39$\times$ & 92.4\% \\
RAG: ML Tutorial & 647 & 5.56$\times$ & 6.46$\times$ & 92.3\% \\
RAG: Python Guide & 676 & 5.51$\times$ & 3.11$\times$ & 92.4\% \\
Long System Prompt & 484 & 3.83$\times$ & 4.74$\times$ & 92.2\% \\
Few-shot Sentiment & 282 & 3.39$\times$ & 3.93$\times$ & 89.4\% \\
\midrule
\textbf{Average} & -- & \textbf{5.13$\times$} & \textbf{5.92$\times$} & 91.8\% \\
\bottomrule
\end{tabular}
\end{table}

\subsubsection{Results on TinyLlama-1.1B}

Table~\ref{tab:targeted_tinyllama} shows results on TinyLlama-1.1B. Speedups are lower than Mistral-7B due to smaller model size, but still achieve \textbf{2.60$\times$} average speedup.

\begin{table}[h]
\centering
\caption{Targeted benchmark results on TinyLlama-1.1B}
\label{tab:targeted_tinyllama}
\begin{tabular}{lcccc}
\toprule
Scenario & Prefix & Speedup & Cached & Reuse \\
\midrule
Code Completion & 1464 & \textbf{3.64$\times$} & \textbf{4.22$\times$} & 93.8\% \\
RAG: API Docs & 660 & 3.94$\times$ & 2.28$\times$ & 94.4\% \\
RAG: ML Tutorial & 656 & 2.96$\times$ & 2.13$\times$ & 94.2\% \\
Long System Prompt & 497 & 1.80$\times$ & 1.83$\times$ & 94.0\% \\
Few-shot Sentiment & 298 & 1.22$\times$ & 1.20$\times$ & 90.8\% \\
\midrule
\textbf{Average} & -- & \textbf{2.60$\times$} & \textbf{2.09$\times$} & 93.4\% \\
\bottomrule
\end{tabular}
\end{table}

\subsubsection{Key Finding: Scenario Matching}

Figure~\ref{fig:scenario_comparison} compares speedups across scenarios. The results reveal that \textbf{scenario matching is crucial}: while multi-turn conversations with growing context show limited gains (0.90$\times$--1.02$\times$), fixed-prefix scenarios achieve 3--10$\times$ speedup.

\begin{figure}[h]
\centering
\includegraphics[width=\columnwidth]{../experiments/figures/scenario_comparison.pdf}
\caption{Speedup comparison across targeted scenarios. Larger models (Mistral-7B) benefit more from caching due to higher per-token computation cost.}
\label{fig:scenario_comparison}
\end{figure}

\subsection{Ablation: Prefix Length Sensitivity}

Figure~\ref{fig:prefix_scaling} shows how speedup scales with prefix length on TinyLlama-1.1B. Short prefixes ($<$200 tokens) show limited benefit, while longer prefixes yield substantial speedups.

\begin{figure}[h]
\centering
\includegraphics[width=0.9\columnwidth]{../experiments/figures/prefix_length_scaling.pdf}
\caption{Speedup scaling with prefix length. Cached speedup exceeds 5$\times$ at 1000 tokens.}
\label{fig:prefix_scaling}
\end{figure}

\begin{table}[h]
\centering
\caption{Prefix length ablation (TinyLlama-1.1B)}
\label{tab:prefix_ablation}
\begin{tabular}{cccc}
\toprule
Prefix & Overall & Cached & Finding \\
\midrule
100 & 0.94$\times$ & 1.14$\times$ & Overhead > benefit \\
250 & 1.98$\times$ & 2.14$\times$ & Breakeven \\
500 & 2.77$\times$ & 3.21$\times$ & Clear benefit \\
750 & 3.12$\times$ & 3.76$\times$ & Strong benefit \\
1000 & \textbf{3.96$\times$} & \textbf{5.05$\times$} & Maximum tested \\
\bottomrule
\end{tabular}
\end{table}

\subsection{Visualization}

Figure~\ref{fig:results} shows a comprehensive visualization of our experimental results.

\begin{figure}[h]
\centering
\includegraphics[width=\columnwidth]{../experiments/results/paper/figures/fig_summary.png}
\caption{Summary of \deltacache{} evaluation: (a) correctness validation showing 100\% token match rate, (b) speedup over baseline, (c) token reuse scaling with prompt length, (d) eviction policy comparison, (e) cache efficiency improvement with query volume.}
\label{fig:results}
\end{figure}

\section{Discussion}
\label{sec:discussion}

\paragraph{Optimal Application Scenarios}
Based on our comprehensive experiments, \deltacache{} is most effective in:
\begin{itemize}
    \item \textbf{Long prefix scenarios ($>$500 tokens)}: RAG with long documents, few-shot learning with many examples, lengthy system prompts
    \item \textbf{High query volume on shared context}: Same document queried multiple times, batch inference with shared prefix
    \item \textbf{TTFT-critical applications}: Interactive systems where first response latency matters
\end{itemize}

\paragraph{When NOT to Use \deltacache{}}
Our real-world dataset validation reveals scenarios where \deltacache{} provides limited or negative benefit:
\begin{itemize}
    \item \textbf{Short prefix ($<$100 tokens)}: Cache overhead exceeds computation savings
    \item \textbf{High new-content ratio}: Multi-turn conversations where each turn adds substantial new content
    \item \textbf{Single-use prompts}: One-off queries with no prefix sharing opportunity
\end{itemize}

\paragraph{Overhead Analysis}
\deltacache{} introduces overhead for prefix tree operations and cache management.
Our experiments show that this overhead becomes amortized when prefix length exceeds $\sim$200 tokens.
For shorter prefixes, the overhead can result in slowdown (0.9$\times$ as observed in ShareGPT-style evaluation).

\paragraph{Correctness and Precision}
We verify 100\% correctness at fp16 precision with tolerance $2 \times 10^{-2}$.
Quantized models (8-bit) show larger numerical differences ($\sim$0.6 max diff) due to quantization noise, not \deltacache{} itself.
For applications requiring strict numerical fidelity, we recommend fp16 or bf16 precision.

\paragraph{Limitations}
Current limitations include:
\begin{itemize}
    \item Single-request focus (no batched inference optimization)
    \item CPU memory offloading adds latency for GPU retrieval
    \item No support for dynamic prefix modification
    \item Overhead not amortized for short prefixes
\end{itemize}

\section{Conclusion}
\label{sec:conclusion}

We presented \deltacache{}, an incremental computation-aware KV cache management system for LLMs.
By leveraging prefix sharing patterns common in production deployments, \deltacache{} eliminates redundant KV computations while maintaining \textbf{100\% numerical correctness} at fp16 precision.

Our comprehensive experiments on TinyLlama-1.1B and Mistral-7B demonstrate:
\begin{itemize}
    \item \textbf{5.13$\times$} average speedup on Mistral-7B (fp16) across targeted application benchmarks
    \item \textbf{10.90$\times$} cached speedup on code completion scenarios (1448 token prefix)
    \item Speedup that scales with prefix length: \textbf{0.94$\times$} (100 tokens) to \textbf{3.96$\times$} (1000 tokens)
    \item \textbf{2.60$\times$} average speedup on TinyLlama-1.1B with up to \textbf{4.22$\times$} on long-prefix workloads
    \item Theoretical complexity reduction from $O(n^2)$ to $O((n-m)^2)$
\end{itemize}

Crucially, our experiments identify \textbf{scenario matching} as the key factor: \deltacache{} excels in fixed-prefix workloads (RAG with shared knowledge base, few-shot prompting, code completion) achieving 3--10$\times$ speedup, while multi-turn conversations with continuously growing context show limited gains (0.90--1.02$\times$).
Ablation studies confirm that speedup is strongly correlated with prefix length, with 500+ tokens being the practical threshold for benefit.

\deltacache{} provides a practical, model-agnostic solution for optimizing LLM inference in scenarios with long shared prompt prefixes.
Target applications include RAG systems querying fixed document collections, few-shot learning with shared examples, code completion with fixed context, and chatbots with lengthy system prompts.
Future work includes extending to batched inference, exploring learned eviction policies, and deeper integration with serving frameworks like vLLM and SGLang.

\paragraph{Code Availability}
\deltacache{} is available as an open-source Python library at \url{https://github.com/anonymous/deltacache}.

\bibliographystyle{plainnat}
\bibliography{references}

\end{document}
